\documentclass[12pt,a4paper]{article}
% ---- Basic packages ----
\usepackage[utf8]{inputenc}
\usepackage[T1]{fontenc}
\usepackage[english]{babel}
\usepackage{lmodern}
% ---- Page layout ----
\usepackage{geometry}
\geometry{a4paper, margin=2.5cm}
\usepackage{enumitem} 
\usepackage[most]{tcolorbox}
    \definecolor{CustomRed}{RGB}{255,51,51}
% ---- Line spacing ----
\usepackage{setspace}
\onehalfspacing  
% ---- Math and symbols ----
\usepackage{amsmath, amssymb, amsfonts}
\usepackage{bm}
\usepackage{siunitx}
\usepackage{listings}
% ---- Tables ----
\usepackage{booktabs, tabularx, longtable, multirow, threeparttable}
% ---- Figures ----
\usepackage{graphicx}
\usepackage{float}
\usepackage{caption}
\usepackage{subcaption}
\usepackage{tabularx}
\usepackage{booktabs}
% ---- Colors and links ----
\usepackage{xcolor}
\usepackage{hyperref}
\usepackage{booktabs}

\setlength{\parindent}{0pt}

\title{Geopolitics for Business \\ Report}

\author{
    Data Group 2 \\[2ex]
    \normalsize Marius Borsu Lefevre \\
    \normalsize Giovanni Clemente \\
    \normalsize Greta Guadalupi \\
    \normalsize Petr Hrebenar \\
    \normalsize Stéphane Miquel Clément \\
    \normalsize Pietro Montano \\
    \normalsize Souhail Habib Nehme \\
    \normalsize Joan Olaf Subirachs Garcia
}

\date{\today}
\begin{document}
\maketitle
\newpage

\section{Literature review}
\subsection{Methods}
Academic and professional notions of the condition of globality, widespread first during the 1980s (Robertson, 1983), emerged along with an increasing cross-border integration of the global economic system resulting after the II World War (Hover, Buchenrieder, and Sokolsky, 2025). Indeed, the characterisation of a liberal regime of financial and trade relations based on the principles of cooperation, mutual interdependence, and openness, prompted the flourishing of multinational enterprises (MNEs) in the international business (IB) environment (Petricevic and Teece 2019; Lessard, Teece, and Leih, 2016). The liberal school of thought, acknowledging the increasing “flatness” of the system (Friedman, 2005; Rugman and Oh, 2008), argued that the resulting deepening of economic integration would diminish the likelihood of inter-state confrontation or institutional crises due to high opportunity costs from the disruption of trade (Mansfield and Pollins, 2003; Choi and Kim, 2021; Gartzke and Li, 2003; Goldsmith, 2013). This political perspective was operationalised, for instance, through Germany’s “change through trade” approach that encouraged the construction of several gas pipelines with the Soviet Union and later Russian Federation between the 1970s and the early 2000s (Blumenau, 2022; Wood, 2023). In light of this, energy infrastructure projects such as Nord Stream 2 were “never just business” (Schmidt-Felzmann, 2016). Geopolitical developments over the past fifteen years have revealed the limitations of the liberal theory, indicating a reversal of an assumption that is reshaping the business landscape once again. 

In this context, several authors have put forward that the current world order is entering a de-globalisation phase (Livesay, 2018; James, 2018); one that is putting at risk, by direct means, the way MNEs operate and do business (Petricevic and Teece, 2019). More specifically, this period of disintegration is making the IB landscape fraught with unprecedented levels of volatility, uncertainty, complexity, and ambiguity (VUCA) (Bennett and Lemoine, 2014; Cavusgil et al., 2021; Schoemaker, Heaton, and Teece, 2018; van Tulder, Verbeke and Jankowska, 2020). As established by Sénquiz-Díaz (2025), the most important contribution to the VUCA framework in the IB context is that of Bennett and Lemoine (2014). For these authors, volatility can be described as a state of “relatively unstable change”, while uncertainty is the “lack of knowledge as to whether an event will have meaningful ramifications” (p. 313). Furthermore, complexity, in the IB environment, reflects the idea that there is no single arrangement, but rather a constellation of systems that interact altogether in an exhaustive “network of information of procedures” (p. 313). Lastly, ambiguity is depicted as the absence of clear guidelines for decision-making, and the margin for independent interpretation in cause-and-effect relationships. 

MNEs are currently operating in such one highly challenging context that not only disrupts routine business operations but also complicates the management of exceptional circumstances. As put forward by Cavusgil et al. (2021), apart from traditional commercial, financial, cultural, and political risks, companies are now dealing with additional and greater threats, including technological obsolescence, cyber security issues, trade wars, sustainability concerns, social challenges, and, more importantly “black swan” events. Indeed, in the systems risk literature, exogenous shocks such as interstate armed conflicts or extreme weather events emerge as “crises” once actual harms are materialised (Homer-Dixon et al., 2015). Nevertheless, as put forward by Lawrence, Janzwood, and Homer-Dixon (2022), today, the causal activation of a given risk in a given system into an actual crisis does not remain confined to the system in which it originates, but it spreads and reinvigorates onto other systems. 

The term polycrisis reflects the confluence and intertwining of these individual crises that “interact, exacerbate, and reshape one another (...) that must be understood and addressed as a whole" (Lawrence, Janzwood, and Homer-Dixon, 2022). The conceptual transition from systemic risk to the notion of polycrisis can be understood through the maximisation of the broader defining properties embedded in the systems that the latter emphasise. Based on the scholarship analysis of Lawrence, Janzwood, and Homer-Dixon (2022), the five distinct characteristics of polycrises are: extreme complexity and dynamism; disproportional cause-effect relationships, transboundary causal links, and profound uncertainty (Renn et al., 2020; 2019; Schweizer and Renn, 2019). Precisely, as put forward by Collste et al. (2025), the identification of valid causality relationships between these crises under such uncertain conditions makes it hard to assess and overcome the challenges derived from it. In the application of a system archetypes approach to identify causal trends within crises, the authors define a model, referred to as “limits to growth”, that effectively specifies our scope of research. In the context of a growing size of economic operations, certain operation constraints and external pressures such as scarce resources or climate-related externalities constrain the expansion capacity of operations of businesses. In this regard, this production slowdown may be already underway in some regions of the world (Collste et al., 2025). Furthermore, in light of this increased complexity, Rapoza (2019) argues that MNEs might be already shifting toward decoupling from their global and regional global value chains (GVCs), a strategy also envisioned by Collste et al. (2025) to manage the implications of polycrises. 

Purposefully, this research answers the call to investigating the causal interactions of polycrisis recently put forward by the scientific community investigating global risk (Future Earth, Sustainability in the Digital Age, and International Science Council, 2021). 
\newpage

\subsection{Metrics}

A growing literature shows that geopolitical shocks affect firms not only through standard financial variables such as investment, financing costs, and profitability, but also through sustainability-related outcomes. In particular, recent firm-level studies find that higher geopolitical risk and geopolitical conflict are associated with weaker ESG performance, including weaker environmental outcomes. This means that environmental performance is not a peripheral variable in this setting, but one of the firm-level margins through which geopolitical stress can be observed [1]–[3].

Once environmental outcomes are accepted as relevant, the next issue is how to measure them. Total capital expenditure (CAPEX) and CSR spending are both informative, but neither is a direct measure of environmental performance. Total CAPEX captures overall investment, but it does not distinguish between investment that reduces environmental harm and investment that maintains or expands emissions-intensive activity. Green CAPEX exists as a reporting concept, but it is still not a universally adopted standalone balance-sheet figure. CSR expenditure is also ambiguous because it may reflect communication, philanthropy, or stakeholder management rather than operational environmental performance. For that reason, an environmental metric tied directly to emissions is closer to the underlying object of interest than either total CAPEX or CSR spending [1].

This is consistent with the broader climate-finance literature, where carbon emissions are treated as a core firm-level climate variable rather than as a secondary disclosure item. Carbon emissions and carbon intensity are regularly used to capture firms’ exposure to climate-related risks and their environmental footprint, and they have been shown to be economically meaningful in studies linking emissions to regulation, firm value, and financial performance [4]–[7]. In that sense, an emissions-based environmental variable is more closely aligned with operational environmental performance than a broader composite sustainability indicator.

Within that class of metrics, Trucost provides environmental data designed to measure firms’ environmental performance quantitatively and systematically. Its methodology defines environmental intensity as environmental impact normalized by revenue, so that the variable reflects environmental efficiency rather than raw firm size alone. This makes the measure comparable across firms and sectors, which is useful when the aim is to study cross-firm differences in environmental performance after a common shock [8].

This is where GHG Direct Intensity fits more specifically. It captures greenhouse-gas emissions generated by the firm’s own operations relative to revenue. In standard emissions accounting, direct emissions correspond to Scope 1 emissions, meaning emissions from sources owned or controlled by the firm. As a result, GHG Direct Intensity measures how emissions-intensive a firm’s own production and operations are, rather than how exposed it is to controversies or how strong its sustainability communication appears to be [8], [9]. This makes it an operational environmental-performance variable.

RepRisk was considered as an alternative environmental metric, but it measures externally reported environmental controversies rather than firms’ underlying operational environmental efficiency. Its methodology is explicitly outside-in, issues- and event-driven, and excludes company self-disclosures, which makes it suitable for tracking incident exposure but less directly tied to emissions intensity or production-related environmental performance. In addition, controversy datasets are affected by media-selection bias: Barkemeyer, Revelli and Douaud (2023) show that the odds of controversy coverage are five times higher for firms headquartered in English-language countries than for firms in other language regions. For that reason, RepRisk is better interpreted as a controversy measure than as a direct measure of firms’ operational environmental performance. [11]

There is also a measurement reason to focus on direct emissions. Recent work by the NGFS shows that inconsistencies across emissions datasets are lowest for direct emissions data, that is, Scope 1, while discrepancies increase for Scope 2 and especially Scope 3. This matters because if only one environmental variable is selected, direct emissions intensity corresponds to the part of firm-level emissions that is both closest to the company’s own activity and less affected by the larger measurement noise found in indirect-emissions categories [10].

The main limitation of Trucost GHG Direct Intensity is frequency. Trucost’s research process is annual rather than quarterly, since its assessments are aligned with firms’ annual reporting cycles [8]. This reduces the ability to trace very short-run post-shock dynamics. At the same time, the main studies linking geopolitical shocks to ESG outcomes also rely on annual firm-level data rather than quarterly panels [1]–[3]. The use of an annual environmental variable is therefore consistent with the frequency that already dominates this literature. The implication is simply that the analysis speaks to medium-run changes in environmental performance after geopolitical shocks, rather than immediate quarter-by-quarter adjustment.

Taken together, the literature implies the following. First, geopolitical shocks are associated with weaker ESG outcomes, including weaker environmental outcomes, so environmental metrics are relevant dependent variables in this setting [1]–[3]. Second, among possible firm-level metrics, a carbon-intensity measure is closer to operational environmental performance than total CAPEX or CSR spending [1], [4]. Third, within carbon metrics, Trucost GHG Direct Intensity measures emissions from the firm’s own operations relative to revenue, which means it captures the emissions intensity of production rather than controversies or self-reported sustainability claims [8], [9]. Its annual frequency is a limitation for short-run timing, but not for its interpretation as a measure of firms’ operational environmental performance [8], [10]. 

\newpage

%\subsection{Dataset construction}

\noindent
This study combines three data layers: a country-level panel, a firm-level Trucost panel, and a final merged regression panel. The objective is to relate changes in the macroeconomic and energy environment to firm-level environmental performance, while accounting for sectoral exposure and country-level conditions. Because firm-level ESG and environmental performance data are not available at quarterly or monthly frequency in a sufficiently consistent manner, the empirical analysis is conducted at the annual level.

\subsubsection{Country panel}

\noindent
The first step consists in constructing a country-level panel for the European countries included in the sample. This panel contains the macroeconomic and energy variables used to characterize the national environment in which firms operate. In particular, it includes electricity prices, inflation, and GDP.

\noindent
Electricity price data are initially collected at daily frequency through an energy market API and then aggregated to the annual level. This aggregation ensures consistency with the frequency of the firm-level environmental data used in the analysis. For each country-year, the resulting measure captures the average electricity price over the year.

\noindent
Macroeconomic controls are added using OECD and World Bank sources. GDP is used to capture the scale of economic activity, while inflation controls for the general price environment. These variables are measured at the country-year level. The resulting country panel therefore provides the macroeconomic and energy context for all firms headquartered in a given country in a given year.

\subsubsection{Trucost firm panel}

\noindent
The second step is the construction of the firm-level environmental panel using Trucost data. The purpose of using Trucost is to rely on a quantitative measure of environmental performance at the firm level rather than on indirect proxies such as CAPEX or broad CSR expenditure.

\noindent
[Placeholder: insert here the literature-based justification for selecting the Trucost variable, in particular why \emph{GHG Direct Intensity} is the preferred proxy for firm-level operational environmental performance in the context of geopolitical or energy shocks.]

\noindent
The firm-level panel is constructed at the annual frequency because ESG and environmental indicators are not available on a quarterly or monthly basis in a sufficiently complete and comparable way. As a result, the empirical design focuses on medium-run variation in environmental performance rather than immediate short-run adjustments.

\noindent
The Trucost sample is then filtered to retain the sectors relevant for the empirical design. Firms are classified into energy-intensive and non-energy-intensive industries in order to define treatment and control groups. Energy-intensive sectors include chemicals, metals and mining, oil and gas, paper and forest products, and construction materials. The control group consists of less energy-intensive but industrially comparable sectors, such as machinery, automobiles, auto components, pharmaceuticals, and biotechnology. Based on this classification, an indicator variable is created that equals one for energy-intensive industries and zero otherwise.

\subsubsection{Final regression panel}

\noindent
In the final step, the firm-level Trucost panel is merged with the country-level panel using country and year identifiers. This produces a firm-country-year dataset in which each observation combines a firm's environmental performance and sectoral classification with the macroeconomic and energy conditions prevailing in its home country.

\noindent
The final regression panel contains three main categories of variables. First, it includes the dependent variable measuring firm-level environmental performance. Second, it includes the treatment dimension based on whether the firm belongs to an energy-intensive industry. Third, it includes country-level covariates such as electricity prices, GDP, and inflation.

\noindent
This merged structure makes it possible to study whether firms in energy-intensive sectors are differentially affected by changes in the energy environment, and whether those changes are associated with variation in firm-level environmental performance. The complete panel therefore links macro-level conditions, sectoral exposure, and firm-level outcomes within a single regression framework.

\section{References}

[1] Abdullah, M., Tiwari, A.K., Hossain, M.R. and Abakah, E.J.A. (2024) ‘Geopolitical risk and firm-level environmental, social and governance (ESG) performance’, \textit{Journal of Environmental Management}, 363, 121245. 

[2] Saharti, M., Chaudhry, S.M., Pekar, V. and Bajoori, E. (2024) ‘Environmental, social and governance (ESG) performance of firms in the era of geopolitical conflicts’, \textit{Journal of Environmental Management}, 351, 119744. 

[3] Erzurumlu, Y.O., Gozgor, G., Lau, C.K.M., Soliman, A.M. and Turkkan, M. (2025) ‘The effects of geopolitical and political risks on corporate ESG practices’, \textit{Journal of Environmental Management}, 384, 125747. 

[4] Sautner, Z., van Lent, L., Vilkov, G. and Zhang, R. (2023) ‘Firm-Level Climate Change Exposure’, \textit{The Journal of Finance}, 78(3), pp. 1449–1498. 

[5] Le, H. and Nguyen-Phung, H. (2024) ‘Assessing the impact of environmental performance on corporate financial performance: A firm-level study of GHG emissions in Africa’, \textit{Sustainable Production and Consumption}, 47, pp. 644–654. 

[6] Perdichizzi, S., Buchetti, B., Cicchiello, A.F. and Dal Maso, L. (2024) ‘Carbon emission and firms’ value: Evidence from Europe’, \textit{Energy Economics}, 131, 107324. 

[7] Bolton, P. and Kacperczyk, M. (2023) ‘Global Pricing of Carbon-Transition Risk’, \textit{The Journal of Finance}, 78(6), pp. 3677–3754.

[8] S\&P Global Sustainable1 (2025) \textit{Trucost Environmental Data: Methodology}.

[9] U.S. Environmental Protection Agency (EPA) (2025) \textit{Scope 1 and Scope 2 Inventory Guidance}.

[10] Network for Greening the Financial System (NGFS) (2024) \textit{Improving Greenhouse Gas Emissions Data}.

[11]Barkemeyer, R., Revelli, C. and Douaud, A. (2023) ‘Selection bias in ESG controversies as a risk for sustainable investors’, \textit{Journal of Cleaner Production}, 405, 137035. https://doi.org/10.1016/j.jclepro.2023.137035. 

Berg, F., Kölbel, J.F. and Rigobon, R. (2022) ‘Aggregate Confusion: The Divergence of ESG Ratings’, \textit{Review of Finance}, 26(6), pp. 1315–1344. \href{https://doi.org/10.1093/rof/rfac033}{https://doi.org/10.1093/rof/rfac033}. 

Bikmetova, N. and Pirinsky, C.A. (2026) ‘Do ESG Rating Agencies Improve ESG Performance?’, \textit{Journal of Business Ethics}, 204, pp. 335–365. https://doi.org/10.1007/s10551-025-06063-0. 

Kathan, M.C., Utz, S., Dorfleitner, G., Eckberg, J. and Chmel, L. (2025) ‘What you see is not what you get: ESG scores and greenwashing risk’, \textit{Finance Research Letters}, 74, 106710. https://doi.org/10.1016/j.frl.2024.106710. 

Marshall, B.R., Nguyen, J.H., Nguyen, N.H., Qiu, B. and Visaltanachoti, N. (2023) \textit{“E” Ratings and Negative Environmental Performance}. Working paper, March 2023 version. 



\end{document}
